\documentclass{jsarticle}
\usepackage{amsmath}
\usepackage{amssymb}
\begin{document}
\title{医学と暦、易学}
\author{Masaaki Yamaguchi}
\date{}
\maketitle

      医学と暦は、占いと自然の摂理に則って作られている。
   しかし、医学では占いは作為体験として分類されている。
   仏教は、気が実在してもいるし、仏眼相にも治癒力としての念力も
   実在している。医学では、占いが作為体験として分類されていないと、
   精神分裂病にとって、平穏な生活ができないために、このような
   見識として見られている。
   ５行は、六星占術が西暦エネルギーと基数エネルギーから
   己の寅の月の人類が誕生した日の十干と干支に、このエネルギーから
   36加群した運命数から、宇宙の十干と干支が本当に数列の和から実在している。
   この数字は六星占術となって、占いが宇宙の十干と干支が出せる数式に
   実際ある。
   墓も治癒力が実在しているが、仏教と医学では、ヒポクラテスの医学から
   食事が大切なために、否定されてもいる。
   六星占術は、占術家から捨てる占いとして認識されている。
   ５行で見ると病気の起因がわかる。医学では、本当に易学で病気を統計から
   見ている。 
   この世は現の世界だから、５行から外れると精神や身体に病気として、
   作為体験となり、病気が現れる。
   作為体験とは、周りからコントロールされている、させられ体験として、
   実体験する。作為体験は、占いから言うと、運命を自身に誘導されていると
   思い込んでいる。占いはストーリでもある。
   だが、医学でも易学は病気を見るのに使われている。自身の性質を見ると、
   ５行のどこの臓器がだめになっているかで、中国医学からは、元型で
   病気を診断する。
    墓は、強力な磁場が墓の周りに発生している。仏教では輪廻転生で350回
    この世に生まれ変わる。
    5行から外れると、自身の性質がこの共鳴している十干と干支からも外れて、
    病気になる。

\end{document}
